\documentclass[14pt]{extarticle}
\usepackage{preamble}
\geometry{letterpaper, landscape, margin=0.5in}
\usepackage{qrcode}
\renewcommand{\answer}[1]{}
\setlength{\parskip}{4pt}

\begin{document}
\pagestyle{empty}

\daybanner{Topic 3.5 \textperiodcentered\ TODAY'S PROBLEM}{Which Design Supports Cause?}

\noindent\begin{minipage}[t]{0.57\linewidth}
\vspace{0pt}
Suppose a researcher recruits 30 student volunteers to study whether an energy drink changes reaction time.\par\smallskip \textbf{Design 1:} The first 15 students to sign up receive the energy drink; the last 15 receive water. A tester who knows each drink measures reaction time.\par \textbf{Design 2:} A coin flip assigns each volunteer to the energy drink or water. The water is flavored to match, both drinks are served in identical opaque cups, and neither the volunteer nor the tester knows which drink was assigned.\par
\end{minipage}\hfill
\begin{minipage}[t]{0.40\linewidth}
\vspace{0pt}
\begin{center}
\textbf{Two energy-drink designs}\par\smallskip
\begin{tabular}{|l|c|c|c|}
\hline
\textbf{\#} & \textbf{Assign} & \textbf{Control} & \textbf{Blind} \\ \hline
\textbf{1} & Signup & Water & No \\ \hline
\textbf{2} & Coin & Matched & Both \\ \hline
\end{tabular}
\end{center}
\end{minipage}

\begin{firsttakebox}{FIRST TAKE (5 min, on your sheet)}
Which design would make you more willing to attribute a reaction-time difference to the drink? Name one feature behind your choice.
\end{firsttakebox}

\begin{description}[leftmargin=1.15in, labelwidth=1in, itemsep=2pt, topsep=2pt]
  \item[\dokbadge{1}~(a)] Identify the experimental units, explanatory variable and its levels, and response variable.
  \item[\dokbadge{2}~(b)] For each design, state whether it has random assignment, a control group, and blinding. Explain one job performed by each feature.
  \item[\dokbadge{3}~(c)] Which design could support the claim that the energy drink caused a reaction-time difference? Cite the specific features that settle it, explain what a difference under Design 1 could be due to instead, and decide whether either result could automatically generalize to all teenagers.
\end{description}

\vfill
\noindent\begin{minipage}{0.84\linewidth}
\textbf{Video + follow-along:} \href{https://robjohncolson.github.io/apstats-live-worksheet/u3_lesson5_live.html}{\texttt{\detokenize{u3_lesson5_live.html}}} (scan the code)\par
\textbf{Finish (a)--(c) after the video. Turn in your sheet.}
\end{minipage}\hfill
\qrcode[height=0.75in]{https://robjohncolson.github.io/apstats-live-worksheet/u3_lesson5_live.html}

\end{document}
